%%%%%%%%%%  scan idx 164  |  folio 151  %%%%%%%%%%
\SGAPrintPage{151}%
\uline{Proof}. We define the arrow by resolving \(F\) injectively and using the canonical arrow of \(\Kpx(S)\)
\[
  E\longrightarrow\underline{\Hom}^\bullet(\underline{\Hom}^\bullet(E,F),F).
\]
To see that the arrow in question is an isomorphism when \(E\) is perfect and \(F=A\), we reduce, by localization and dévissage, to \(E=A\), in which case it is evident.

\sgaStableTarget{sga6.e01.def.7-3}\textbf{Definition 7.3.} \uline{Let} \(E\in\Ob\Dpx(S)\). \uline{We shall put} \(\mathbb{R}\underline{\Hom}(E,A)=E^\vee\), \uline{and say that} \(E^\vee\) \uline{is the} dual \uline{of} \(E\).

By (7.1) and (7.2), when \(E\in\Ob\Dpx(S)\) is perfect (resp.\ perfect and of finite tor-dimension), \(E^\vee\) is also, and the canonical arrow \(E\to E^{\vee\vee}\) is an isomorphism. In other words, the functor \(\mathbb{R}\underline{\Hom}(\quad,A)\) defines an anti-auto-equivalence of the category \(\Dpx(S)_{\mathrm{parf}}\) (resp.\ \(\Dpx(S)_{\underline{\mathrm{parf}}}\) (5.8.4.b))) with itself.

\sgaStableTarget{sga6.e01.lem.7-4}\textbf{Lemma 7.4} (CARTAN isomorphism) (cf.\ ([H] II 5.15)). \uline{For} \(E\in\Ob\Dpx(S)\), \(F\in\Ob\Dpx^-(S)\), \(G\in\Ob\Dpx^+(S)\), \uline{there exist canonical functorial isomorphisms}:
\[
\begin{aligned}
\mathrm R\underline{\Hom}(E\otL F,G)&\xrightarrow{\sim}\mathrm R\underline{\Hom}(E,\mathrm R\underline{\Hom}(F,G)),\\
\RHom(E\otL F,G)&\xrightarrow{\sim}\RHom(E,\mathrm R\underline{\Hom}(F,G)),\\
\Hom(E\otL F,G)&\xrightarrow{\sim}\Hom(E,\mathrm R\underline{\Hom}(F,G)).
\end{aligned}
\]

%%%%%%%%%%  scan idx 165  |  folio 152  %%%%%%%%%%
\SGAPrintPage{152}%
\uline{Proof}: This is standard: resolve \(G\) injectively and \(F\) flatly.

\sgaStableTarget{sga6.e01.lem.7-5}\textbf{Lemma 7.5.} \uline{For} \(E\in\Ob\Dpx^-(S)\), \(F\in\Ob\Dpx^+(S)\), \uline{there exists a canonical functorial homomorphism}
\[
  E\otL\mathrm R\underline{\Hom}(E,F)\longrightarrow F
\]

\sgaStableTarget{sga6.e01.prop.7-6}\textbf{Proposition 7.6} (cf.\ ([H] II 5.13)). \uline{Let} \(E\in\Ob\Dpx(S)\), \(F\in\Ob\Dpx^+(S)\), \(G\in\Ob\Dpx^+(S)\). \uline{Under one of the following hypotheses}
\begin{enumerate}[label=\roman*)]
\item \(G\) \uline{of finite tor-dimension},
\item \(E\in\Ob\Dpx^-(S)\), \(F\) \uline{of finite tor-dimension and} \(G\in\Ob\Dpx^b(S)\),
\end{enumerate}
\uline{there exists a canonical functorial homomorphism}
\[
  \mathrm R\underline{\Hom}(E,F)\otL G
  \longrightarrow
  \mathrm R\underline{\Hom}(E,F\otL G).
\]
\uline{It is an isomorphism when} \(E\) \uline{is perfect}.

\uline{Proof}. Let us define the arrow. Under hypothesis (i), it suffices to resolve \(F\) injectively and \(G\) flatly. Under hypothesis (ii), it is equivalent, by (7.4), to define an arrow
\[
  E\otL\mathbb{R}\underline{\Hom}(E,F)\otL G\longrightarrow F\otL G.
\]
We apply the functor \(\otL G\) to the arrow of (7.5).

To see that the arrow is an isomorphism when \(E\) is perfect, we use the standard method (cf.\ the proof of (7.2)): we reduce, by localization and dévissage, to the case \(E=A\).

%%%%%%%%%%  scan idx 166  |  folio 153  %%%%%%%%%%
\SGAPrintPage{153}%
\sgaStableTarget{sga6.e01.cor.7-7}\textbf{Corollary 7.7.} \uline{There exists, for} \(E\in\Ob\Dpx(S)_{\underline{\mathrm{parf}}}\) and \(F\in\Ob\Dpx^+(S)\), \uline{a functorial isomorphism}
\[
  E^\vee\otL F\xrightarrow{\sim}\mathbb{R}\underline{\Hom}(E,F),
\]
\uline{with the notation of} (7.3).

\uline{Proof}. Apply (7.6), replacing \(F\) and \(G\) respectively by \(A\) and \(F\), and use the fact (7.1) that \(E^\vee\in\Ob\Dpx(S)_{\underline{\mathrm{parf}}}\).

\sgaStableTarget{sga6.e01.cor.7-8}\textbf{Corollary 7.8.} \uline{Let} \(E,F\in\Ob\Dpx(S)_{\underline{\mathrm{parf}}}\). \uline{There exists a natural arrow}
\[
  \mathbb{R}\underline{\Hom}(E,F)\otL\mathbb{R}\underline{\Hom}(F,E)\longrightarrow A,
\]
\uline{which is a} ``\uline{perfect duality}'', i.e.\ \uline{identifies, modulo} (7.4), \uline{each of the complexes} \(\mathbb{R}\underline{\Hom}(E,F)\), \(\mathbb{R}\underline{\Hom}(F,E)\) \uline{with the dual of the other} (\uline{in the sense of} (7.3)).

\uline{Proof}. By (7.5), we have natural pairings \(E^\vee\otL E\to A\), \(F^\vee\otL F\to A\), hence, by tensor product, a natural pairing
\[
  E^\vee\otL F\otL F^\vee\otL E\longrightarrow A,
\]
which, in view of (7.7), gives the stated pairing on the \(\mathbb{R}\underline{\Hom}\). Let us show that it is a perfect duality. We have the canonical isomorphisms
\[
\begin{aligned}
\mathbb{R}\underline{\Hom}(F^\vee\otL E,A)&\simeq\mathbb{R}\underline{\Hom}(F^\vee,E^\vee) &&(7.4)\\
&\simeq F^{\vee\vee}\otL E^\vee &&(7.6)\\
&\simeq F\otL E^\vee &&(7.2).
\end{aligned}
\]
Thus, after checking compatibility, we are done. We could also localize and use dévissage to reduce to \(E=F=A\).

%%%%%%%%%%  scan idx 167  |  folio 154  %%%%%%%%%%
\SGAPrintPage{154}%
\sgaStableTarget{sga6.e01.exercise.7-9}\textbf{Exercise 7.9.} Let \(E,F\in\Ob\Dpx(S)_{\underline{\mathrm{parf}}}\). Show that there exists a canonical isomorphism
\[
  E^\vee\otL F^\vee \xrightarrow{\sim} (E\otL F)^\vee .
\]

\sgaStableTarget{sga6.e01.sec.8}\section*{8. Traces and cup-products}

\sgaStableTarget{sga6.e01.para.8-0}8.0. In this section, we generalize to perfect complexes the notion of the trace of an endomorphism of a projective module of finite type (see SGA 5 III for applications). The hypotheses and notation are those of (7.0).

\sgaStableTarget{sga6.e01.para.8-1}8.1. Let \(E\in\Ob\Dpx(S)_{\underline{\mathrm{parf}}}\) (5.8.4.b)). There exists a natural arrow
\[
  \tau:\mathbb{R}\underline{\Hom}(E,E)\longrightarrow A,
\]
composed of the inverse isomorphism of (7.7), \(\mathbb{R}\underline{\Hom}(E,E)\xrightarrow{\sim}E^\vee\otL E\), and the canonical arrow (7.5) \(E^\vee\otL E\to A\).

Let \(U\in\Ob \underline{S}\). Applying the functor \(H^0(U,.)\) to \(\tau\), and using the fact that \(H^0(U,\mathbb{R}\underline{\Hom}(E,E))\) is identified with \(\Hom_{\Dpx(U)}(E|U,E|U)\), we deduce a homomorphism of \(H^0(U,A)\)-modules,
\[
  \Tr_{E|U}:\Hom_{\Dpx(U)}(E|U,E|U)\longrightarrow H^0(U,A),
\]
which is called the \uline{trace homomorphism}\SGAAuthorFootnote{raisedstar}{The finite tor-dimension hypothesis on \(E\) is in fact superfluous for the definition of \(\Tr_{E|U}\), see (8.5).}.

Before going further, let us give some properties of the trace that follow immediately from the definition.

%%%%%%%%%%  scan idx 168  |  folio 155  %%%%%%%%%%
\SGAPrintPage{155}%
\sgaStableTarget{sga6.e01.para.8-1-1}8.1.1. When \(E=A\), \(\Tr_{A|U}:\Hom(A|U,A|U)\longrightarrow H^0(U,A)\) is precisely the well-known isomorphism. When \(E\) is a direct factor of a free Module of finite type, a homomorphism from \(E\) to \(E\) over \(U\) is locally defined by
\[
  e\longmapsto \sum\langle e,x_i^\vee\rangle x_i,
\]
where \(x_i\) (resp.\ \(x_i^\vee\)) is a local section of \(E\) (resp.\ \(E^\vee\)); the corresponding trace is \(\sum\langle x_i,x_i^\vee\rangle\). We easily deduce that, when \(E\) and \(F\) are direct factors of free Modules of finite type, and \(u:(E+F)|U\longrightarrow(E+F)|U\) is defined by a matrix \(\begin{pmatrix}a&b\\ c&d\end{pmatrix}\), we have
\[
  \Tr_{(E+F)|U}(u)=\Tr_{E|U}(a)+\Tr_{F|U}(d).
\]

\sgaStableTarget{sga6.e01.para.8-1-2}8.1.2. Let \(E\) be a strictly perfect complex, \(U\in\Ob \underline{S}\), and \(u:E|U\longrightarrow E|U\) a homomorphism of complexes, with components \(u^i:E^i|U\to E^i|U\). We then have
\[
  \Tr_{E|U}(u)=\sum(-1)^i\Tr_{E^i|U}(u^i).
\]
Indeed, since the question is local, we can suppose that \(u\) is defined by a section of \(\underline{H}^0(E^\vee\otimes E)\), and even that this section is of the form \(\coprod_i\sum_j x_{ij}^\vee\otimes x_{ij}\), where \(x_{ij}\) (resp.\ \(x_{ij}^\vee\)) is a section of \(E^i\) (resp.\ \(E^{i\vee}\)). Making the arrow (7.5) \(E^\vee\otimes E\to A\) explicit, we find
\[
  \Tr_{E|U}(u)=\sum(-1)^i\langle x_{ij},x_{ij}^\vee\rangle
  =\sum(-1)^i\Tr_{E^i|U}(u^i),
\]
q.e.d.

\sgaStableTarget{sga6.e01.para.8-2}8.2. Let \(E,F\in\Ob\Dpx(S)_{\underline{\mathrm{parf}}}\). We have seen (7.8) that there exists a canonical homomorphism
\[
  \gamma:\mathbb{R}\underline{\Hom}(E,F)\otL\mathbb{R}\underline{\Hom}(F,E)\longrightarrow A.
\]
Let \(U\in\Ob \underline{S}\), and take \(u\in H^0(U,\mathbb{R}\underline{\Hom}(E,F))\) and \(v\in H^0(U,\mathbb{R}\underline{\Hom}(F,E))\).

